\clearpage
\part*{Operational and geometric extensions}
\addcontentsline{toc}{part}{Operational and geometric extensions}

\section{Prescribed success probability: an exact two-ball geometry}\label{sec:lens}
\label{sec:probability}

The unrestricted quantity ignores how frequently the chosen outcome occurs.
Filtering-probability tradeoffs have been studied for particular systems \cite{probabilityprior}. We now obtain an arbitrary-state one-sided formula. For $0<p\leq1$ define
\[
 \HF_A(\rho;p)=\sup\left\{F_3(\rho_K):K^\dagger K\preceq I,
     \ \tr[(K^\dagger K\otimes I)\rho]\geq p\right\},
 \qquad V(p)=\HF_A(\rho;p)^2.
\]
The supremum may be restricted to invertible filters without changing its
value. We prove that assertion rather than assume attainment. All the
results in this part have a Bob version by the same swap as in
\eqref{eq:swap}. We retain the canonical notation
\[
 C=SS^T,\quad g=Sc,\quad d=\tr C-\norm c^2,\quad H=C-dI,
 \quad q(x)=\tau+2g^Tx+x^THx.
\]

\begin{theorem}[Exact probability lens]\label{thm:lens}
If $\norm a<1$, then
\begin{equation}\label{eq:lens}\boxed{
 V(p)=\max_{x\in L_p}q(x),\qquad
 L_p=\{x:\norm x\leq1,\ \norm{a-px}\leq1-p\}.}
\end{equation}
For $0<p<1$ this is the intersection of the two explicit balls
$B(0,1)$ and $B(a/p,(1-p)/p)$. Each feasible point is realized with
success probability exactly $p$ by
\begin{equation}\label{eq:ratefilter}
 K=\sqrt p\,U\sqrt{I+x\cdot\sigma}\,W.
\end{equation}
At $p=1$, $L_p=\{a\}$ and $V(1)=F_3(\rho)^2$.
If $\norm a=1$, the entire frontier is instead constant:
$\HF_A(\rho;p)=\norm b$ for every $0<p\leq1$.
\end{theorem}
\begin{proof}
An admissible filter of probability $P\geq p$ can be multiplied by
$\sqrt{p/P}$ without changing its output. Thus it suffices to prescribe
probability exactly $p$. The canonical marginal is $I/2$, so
\eqref{eq:ratefilter} has that probability. Positivity of its canonical
effect is exactly $\norm x\leq1$. Moreover
\[
 K^\dagger K\preceq I
 \ \Longleftrightarrow\ p(I+x\cdot\sigma)\preceq I+a\cdot\sigma
 \ \Longleftrightarrow\ \norm{a-px}\leq1-p.
\]
Polar decomposition proves that this parametrization covers every output.
The lens is compact and nonempty, since it contains $a$.

For $p<1$, $a$ is strictly inside both balls: $\norm a<1$ and
$\norm{a-pa}=(1-p)\norm a<1-p$. Consequently
$x_\varepsilon=(1-\varepsilon)x+\varepsilon a$ remains feasible and has
$\norm{x_\varepsilon}<1$. Continuity of $q$ shows that the closed-lens
maximum equals the invertible-filter supremum. At $p=1$, the only point
is $a$, and $\sqrt{I+a\cdot\sigma}W=I$. The pure-marginal statement
follows from the product-state reduction already proved.
\end{proof}

\begin{corollary}[Containment regimes and a universal achievable rate]
\label{cor:regimes}
Put $r_a=\norm a$, $p_-=(1-r_a)/2$ and $p_+=(1+r_a)/2$.
For $p\leq p_-$ the lens is the full unit ball. For $p_+\leq p<1$
it is the second ball. Only $p_-<p<p_+$ requires a genuine lens.
Every unrestricted optimum in the closed filter domain therefore has
some realization with success at least $p_-$. Every strictly violating
optimum has a realization with success strictly greater than $p_-$.
\end{corollary}
\begin{proof}
The two containments are respectively equivalent to
$r_a/p+1\leq(1-p)/p$ and
$r_a/p+(1-p)/p\leq1$. For the strict assertion use the unique interior
optimizer and
\[
 W(I+x\cdot\sigma)W\preceq
   \frac{1+\norm x}{1-\norm a}I,
 \qquad p_{\max}(x)\geq\frac{1-\norm a}{1+\norm x}>p_-.
\]
\end{proof}

\begin{corollary}[Canonical frontier: a single six-dimensional spectrum]
\label{cor:canonicalfrontier}
If $a=0$ and $0<p<1$, put $R=\min\{1,(1-p)/p\}>0$. Define
\[
 \nu_R=\rmax\begin{pmatrix}Q&I\\Qcc^T/R^2&Q\end{pmatrix},
 \qquad\eta_R=\max\{d,\nu_R\}.
\]
Then
\begin{equation}\label{eq:canonicalfrontier}\boxed{
 V(p)=\tau+R^2(\eta_R-d)+c^TQ(\eta_RI-Q)^\dagger c.}
\end{equation}
In particular the unrestricted canonical optimum is achievable with
probability at least $1/2$, and a strictly violating optimum with
probability strictly greater than $1/2$.
\end{corollary}
\begin{proof}
The two balls are concentric. Apply the quadratic theorem to $x=Ry$.
Equivalently, the regular secular equation is
$\norm{(\eta_RI-C)^{-1}g}^2=R^2$, with the same kernel completion in
the hard case. For $x_0=(\eta_RI-C)^\dagger g$ the exact identity is
\[
 V(p)-q(x)=(\eta_R-d)(R^2-\norm x^2)
       +(x-x_0)^T(\eta_RI-C)(x-x_0).
\]
All maximizers follow by setting its nonnegative terms to zero.
\end{proof}

For completeness, the translated-ball branch is equally explicit. On
$B(u,R)$, $R>0$, set $g_u=g+Hu$, $\tau_u=q(u)$ and
\[
 \nu_u=\rmax\begin{pmatrix}H&I\\g_ug_u^T/R^2&H\end{pmatrix},
 \qquad\mu_u=\max\{0,\nu_u\}.
\]
Its maximum is
\begin{equation}\label{eq:translatedball}
 \tau_u+\mu_uR^2+g_u^T(\mu_uI-H)^\dagger g_u.
\end{equation}
Here $u=a/p$, $R=(1-p)/p$. A zero radius is the single-point case.

\section{A finite algebraic formula for the genuine lens}
\label{sec:finite_lens}

Generic quadratic optimization over two balls is established mathematics
\cite{twoballs,rontsis}. The contribution of this section is the explicit
fixed-dimensional candidate formula for the physical domain
\eqref{eq:lens}, including feasible singular families and all relevant
stationary roots. It does not assume that a two-constraint S-lemma is exact.

\begin{theorem}[At most fourteen candidate values]\label{thm:fourteen}
Let $q(x)=\tau+2g^Tx+x^THx$ with $H=H^T\in\R^{3\times3}$.
In the transverse lens regime of Corollary \ref{cor:regimes}, its maximum
is the largest of the following explicitly constructed feasible values:
one stationary-affine value, at most six stationary values from each
boundary sphere, and one common-circle value. Every feasibility test below
uses only arithmetic, square roots, matrix spectra and pseudoinverses.
There is no remaining continuous optimization.
\end{theorem}

\paragraph{Candidate 1: the stationary affine family.}
Write the two ball centres and radii as $(u_0,R_0)$ and $(u_1,R_1)$.
For \eqref{eq:lens} these are $(0,1)$ and $(a/p,(1-p)/p)$.
If $g\notin\Ran H$, there is no candidate of this type. Otherwise set
\[
 x_s=-H^\dagger g,\qquad K=\ker H,
 \qquad V_s=\tau-g^TH^\dagger g,
\]
and let $P_K$ be orthogonal projection onto $K$. Define
\[
 k_i=P_K(u_i-x_s),\qquad
 s_i^2=R_i^2-\norm{(I-P_K)(x_s-u_i)}^2.
\]
Include $V_s$ exactly when $s_0^2,s_1^2\geq0$ and
\begin{equation}\label{eq:affinefeasible}
 \norm{k_0-k_1}\leq\sqrt{s_0^2}+\sqrt{s_1^2}.
\end{equation}
A witness is $x_s+z$ with $z=k_0$ if $k_0=k_1$; otherwise use
\[
 z=k_0+\min\left\{1,\frac{\sqrt{s_0^2}}{\norm{k_1-k_0}}\right\}(k_1-k_0).
\]
This remains valid when $K=\{0\}$.

\paragraph{Candidates 2--13: all stationary sphere families.}
For $i\in\{0,1\}$ put $g_i=g+Hu_i$ and $\tau_i=q(u_i)$.
Choose an orthonormal eigenbasis of $H$ with eigenvalues $h_j$ and
components $g_{ij}$ of $g_i$. Form the degree-six polynomial
\begin{equation}\label{eq:lenspoly}
 P_i(\ell)=R_i^2\prod_{j=1}^3(\ell-h_j)^2
    -\sum_{j=1}^3g_{ij}^2\prod_{k\ne j}(\ell-h_k)^2.
\end{equation}
The roots also are the eigenvalues of
$\left(\begin{smallmatrix}H&I\\g_ig_i^T/R_i^2&H\end{smallmatrix}\right)$.
Every \emph{real} root must be considered, including roots below the
largest real root.

If $\ell\notin\operatorname{spec}H$, set
\[
 y=(\ell I-H)^{-1}g_i,\qquad
 V_i(\ell)=\tau_i+\ell R_i^2+g_i^Ty.
\]
Include this value exactly when $\norm{u_i+y-u_j}\leq R_j$, $j\ne i$.
The root equation already gives $\norm y=R_i$.

For a spectral value $\ell$, put $K_\ell=\ker(\ell I-H)$.
Require $g_i\perp K_\ell$ and
$y_0=(\ell I-H)^\dagger g_i$ with $\norm{y_0}\leq R_i$.
Set $s=\sqrt{R_i^2-\norm{y_0}^2}$ and let $P_\ell$ project onto $K_\ell$.
The family and its constant value are
\[
 x=u_i+y_0+z,\quad z\in K_\ell,\quad\norm z=s,
 \qquad V_i(\ell)=\tau_i+\ell R_i^2+g_i^Ty_0.
\]
For $w=u_i+y_0-u_j$, include the value exactly when
\begin{equation}\label{eq:spherefeasible}
 \norm{(I-P_\ell)w}^2+(\norm{P_\ell w}-s)^2\leq R_j^2.
\end{equation}
One witness chooses $z=-sP_\ell w/\norm{P_\ell w}$ when the denominator
is nonzero; otherwise any length-$s$ kernel vector works. Every compatible
singular multiplier is also a root of \eqref{eq:lenspoly}, so these
families add no values beyond its six possible distinct real roots.

\paragraph{Candidate 14: the common circle.}
For the physical lens define
\[
 h=1-\frac{1-\norm a^2}{2p},\qquad
 z=\frac{h}{\norm a^2}a,\qquad
 r=\sqrt{1-\frac{h^2}{\norm a^2}}>0.
\]
Choose $U\in\R^{3\times2}$ with orthonormal columns spanning $a^\perp$,
and put
\[
 H_2=U^THU,\qquad g_2=U^T(g+Hz),\qquad\tau_2=q(z),
 \qquad\nu_2=\rmax\begin{pmatrix}H_2&I_2\\g_2g_2^T/r^2&H_2\end{pmatrix}.
\]
Always include
\begin{equation}\label{eq:circlevalue}
 V_{\rm circ}=\tau_2+\nu_2r^2+g_2^T(\nu_2I_2-H_2)^\dagger g_2.
\end{equation}
There is \emph{no} clamp by zero in this expression: a sphere equality
allows a negative multiplier. Its maximizing points are
$z+U(y_0+v)$ where $y_0=(\nu_2I_2-H_2)^\dagger g_2$,
$v\in\ker(\nu_2I_2-H_2)$ and $\norm{y_0+v}=r$.

\begin{proof}[Proof of Theorem \ref{thm:fourteen}]
The stationary equation inside both balls is $Hx=-g$. Its complete
solution set is $x_s+K$, on which $q=V_s$. Orthogonal projection turns
the two original ball constraints into the two balls
$B_K(k_i,\sqrt{s_i^2})$ in $K$. They meet exactly under
\eqref{eq:affinefeasible}; the displayed witness proves sufficiency.

On a positive-radius sphere the constraint gradient never vanishes.
Lagrange multipliers therefore give
$(\ell I-H)(x-u_i)=g_i$. Off the spectrum, the norm equation is precisely
\eqref{eq:lenspoly}. At a spectral value, orthogonality to the kernel is
necessary and sufficient for solvability; the pseudoinverse solution is
orthogonal to the kernel. Substitution into the quadratic gives the stated
constant family value. Minimizing distance to the other centre over this
kernel sphere is elementary: decompose $w$ orthogonally and choose the
kernel vector opposite $P_\ell w$. The resulting squared distance is
exactly the left side of \eqref{eq:spherefeasible}. A compatible singular
multiplier is a root because a nonzero vector in its kernel is killed by
$(\ell I-H)^2-g_ig_i^T/R_i^2$.

Subtracting the two sphere equations gives $a\cdot x=h$. Their
intersection is the stated circle. Restriction to its affine plane yields
the two-dimensional quadratic. The greatest-real-root proof from the
main theorem applies to a sphere without imposing a nonnegative
multiplier: $\nu_2\geq\lambda_{\max}(H_2)$ still makes its matrix gap
positive semidefinite, and the radial term vanishes identically.
This proves \eqref{eq:circlevalue} and its maximizers, including hard cases.

Finally the lens is compact. A global maximizer is either inside both
balls, on exactly one sphere, or on their common circle. The first case
is stationary; the second is a local extremum on that sphere; the third
is a global extremum on the circle. The three constructions exhaust
these possibilities. Conversely every included value has an explicit
feasible witness. The largest candidate is therefore exactly the global
maximum. All maximizers are the feasible members of precisely those
families whose scalar value equals that maximum. The count is
$1+6+6+1=14$.
\end{proof}

\begin{proposition}[A physical nonglobal sphere maximum is indispensable]
\label{prop:nonglobal}
Take the full-rank physical Bloch data
\[
 a=(-15/17,0,0),\quad b=(-14/85,0,0),\quad
 T=\diag(18/85,2/85,4/425),\qquad p=10/51.
\]
The unique lens optimizer is $x_*=(-2/5,0,0)$ and
\begin{equation}\label{eq:nonglobalvalue}
 V(p)=\frac{17509}{250000}.
\end{equation}
It is a strict local maximum on the second boundary sphere, but not that
sphere's global maximum. Keeping only the largest secular root on each
sphere and the common circle gives an incorrect answer.
\end{proposition}
\begin{proof}
The canonical data are $c=(1/10,0,0)$ and
$S=\diag(3/10,1/20,1/50)$. They are strictly positive since
$\norm c+\sum_i|S_{ii}|=47/100<1$. Inverse whitening with the stated
$a$ gives the displayed rational physical data, using
$\sqrt{1-\norm a^2}=8/17$. The second ball has centre $(-9/2,0,0)$
and radius $41/10$, and
\[
 q(x)=\frac{929}{10000}+\frac3{50}x_1+\frac{71}{10000}x_1^2
       -\frac{201}{2500}x_2^2-\frac{33}{400}x_3^2.
\]
Every feasible $x$ has $-1\leq x_1\leq-2/5$. Setting $x_2=x_3=0$
increases $q$, and the remaining function has derivative at least
$229/5000>0$ on that interval. This proves uniqueness and
\eqref{eq:nonglobalvalue}.

On the second sphere the translated gradient is
$g_1=(-39/20000,0,0)$ and the multiplier at $x_*$ is
$\ell=679/102500>0$. Its tangent Hessian is negative definite.
The opposite pole $(-43/5,0,0)$ gives the global sphere maximum and
is infeasible for the unit ball. The unit-sphere global maximum
$(1,0,0)$ is infeasible for the second ball. The unconstrained stationary
point has $x_1=-300/71$, and the common-circle maximum is only
$118/28125<17509/250000$. All claims are also checked with exact
rational arithmetic in the executable verification.
\end{proof}

\section{The highest success rate preserving the unrestricted optimum}
\label{sec:bestrate}

The rate of an arbitrarily chosen optimal filter need not be the best rate
within a degenerate optimal family. The following formula solves that
secondary problem without optimization.

\begin{theorem}[Best rate and exact plateau endpoint]\label{thm:bestrate}
Assume $\norm a<1$. Let $A=\eta I-C$, $x_0=A^\dagger g$, $\lambda=\eta-d$ as in the main
theorem. Put
\[
 K=\ker A,\quad a_K=P_Ka,\quad a_R=a-a_K,\quad
 r=\sqrt{1-\norm{x_0}^2},\quad\alpha=\norm{a_K},\quad u=a\cdot x_0,
\]
and define
\begin{equation}\label{eq:pb}
 p_b=\frac{1-\norm a^2}{2(1-u-r\alpha)}.
\end{equation}
If $\lambda>0$ or $r=0$, the greatest rate of an unrestricted optimal
filter in the closed filter domain is $p_{\rm best}=p_b$.
If $\lambda=0$, $r>0$ and $p_br\leq\alpha$, the same formula holds.
In all other cases,
\begin{equation}\label{eq:pi}\boxed{
 p_{\rm best}=p_i=
 \frac{1-\norm{a_R}^2}
 {1-u+\sqrt{(1-u)^2-r^2(1-\norm{a_R}^2)}}.}
\end{equation}
For the boundary branch choose
$x=x_0+r a_K/\alpha$ when $\alpha>0$; any available length-$r$
kernel direction works when $\alpha=0$.
For the interior branch choose $x=x_0+a_K/p_i$.
Moreover the entire probability frontier has the exact plateau
\begin{equation}\label{eq:plateau}
 V(p)=\HF_A(\rho)^2\quad\Longleftrightarrow\quad0<p\leq p_{\rm best}.
\end{equation}
If the filtered marginal is pure instead, the frontier is constant and
$p_{\rm best}=1$.
\end{theorem}
\begin{proof}
The unrestricted optimizing family is $x_0+z$ with $z\in K$,
$\norm z\leq r$, and equality required if $\lambda>0$.
On its boundary the rate \eqref{eq:success} becomes
$\delta/[2(1-a\cdot x)]$. Maximizing the linear functional $a\cdot x$
on this kernel sphere gives $u+r\alpha$, proving \eqref{eq:pb}.
The denominator is strictly positive because $u+r\alpha\leq\norm a<1$.

When the full disk is permitted, probability $p$ is feasible on it exactly
when the minimum squared distance to $a$ obeys
\begin{equation}\label{eq:diskdistance}
 \norm{a_R-px_0}^2+\max\{\alpha-pr,0\}^2\leq(1-p)^2.
\end{equation}
This formula is the orthogonal decomposition of the distance to the
radius-$pr$ kernel ball. For $pr\leq\alpha$, expansion gives precisely
$p\leq p_b$. If $p_br<\alpha$, the first infeasible probabilities
therefore occur immediately beyond $p_b$. At the threshold $p_br=\alpha$,
the left side of \eqref{eq:diskdistance} minus its right side has right
derivative $2(1-u-r\alpha)>0$ at $p_b$, so infeasibility likewise begins
immediately after $p_b$. Feasible probabilities form an
interval, since a physical filter can always be scaled down. Thus no later
feasible interval can occur.

If $p_br>\alpha$, feasibility reaches $pr=\alpha$. For $pr\geq\alpha$,
\eqref{eq:diskdistance} becomes
\[
 -r^2p^2+2(1-u)p-(1-\norm{a_R}^2)\leq0.
\]
Its smaller positive root, after rationalization, is \eqref{eq:pi}.
The larger root is at least one, since at $p=1$ the left side equals
$\norm{a_R-x_0}^2\geq0$. If the larger root equals one, then $a_R=x_0$
and both roots equal one, so no isolated feasible endpoint is omitted.
The smaller root is consequently the last
feasible probability in $[0,1]$. It lies beyond $\alpha/r$, so the stated
kernel correction is interior and attains the rate.
The radicand is nonnegative: Cauchy--Schwarz bounds it below by
$(\norm{a_R}-\norm{x_0})^2$.

Finally $V(p)$ equals the unrestricted value precisely when the compact
lens intersects the compact unrestricted maximizing family. The largest
rate for such an intersection is exactly $p_{\rm best}$, proving
\eqref{eq:plateau} and strict loss above it.
For a pure filtered marginal, the product-state reduction makes every
successful filtered value the same, and the identity attains it at rate one.
\end{proof}

\paragraph{The invertible convention.}
If no invertible filter attains the unrestricted $F_3$ value, the set of
exact invertible optimizers is empty; a supremum over that empty set must
not be substituted for $p_{\rm best}$. In this case $p_{\rm best}$ is the
greatest limiting success rate of \emph{asymptotically} optimal invertible
filters. To prove this, approach a best-rate boundary point from the ball
interior; conversely any sequence of asymptotically optimal effects has
a subsequence converging to the closed optimal family, and
$p_{\max}(x)$ is continuous. If $\lambda=0$ and $r>0$, invertible exact
optimizers exist, and the formula is their rate supremum, attained
exactly when some best-rate point is interior. Strict violation always
falls in the attained interior case.

\section{Continuity of the probability frontier}

\begin{proposition}[Monotonicity and an explicit modulus]\label{prop:ratecontinuity}
For a mixed filtered marginal, $V(p)$ is nonincreasing and continuous on
$(0,1]$. For $0<p<p'\leq1$,
\begin{equation}\label{eq:ratemodulus}
 0\leq V(p)-V(p')\leq
 2(\norm g+\norm H_{\rm op})(1+\norm a)\frac{p'-p}{pp'}.
\end{equation}
It is constant near zero, with the exact plateau endpoint from
Theorem \ref{thm:bestrate}.
\end{proposition}
\begin{proof}
Feasibility is preserved when a filter is scaled down, so $L_{p'}\subseteq L_p$.
For $x\in L_p$, let
\[
 s=\frac{p(1-p')}{p'(1-p)},\qquad x'=sx+(1-s)a.
\]
Then $0\leq s\leq1$, $\norm{x'}\leq1$, and
$p'x'-a=[(1-p')/(1-p)](px-a)$, hence $x'\in L_{p'}$.
Also
\[
 \norm{x-a}\leq\frac{(1-p)(1+\norm a)}p,
 \qquad\norm{x'-x}\leq(1+\norm a)\frac{p'-p}{pp'}.
\]
On the unit ball $\norm{\nabla q}\leq2(\norm g+\norm H_{\rm op})$.
Apply this Lipschitz bound to a maximizing $x\in L_p$ and its feasible
$x'$, and use monotonicity. The endpoint $p'=1$ is included.
\end{proof}

\begin{theorem}[A fixed positive rate removes every state discontinuity]
\label{thm:fixedratecontinuous}
For every fixed $p\in(0,1]$, the map
$\rho\mapsto\HF_A(\rho;p)$ is continuous on the entire compact set of
two-qubit density matrices, including all pure-marginal states.
The same holds for Bob.
\end{theorem}
\begin{proof}
Use all complex contractions $K$, a compact set in finite dimension.
The admissibility condition $\tr(K\rho_AK^\dagger)\geq p$ is closed, and
the normalized output is continuous whenever that denominator is at least
$p>0$. If $\rho_n\to\rho$, any subsequential limit of maximizing
contractions is feasible for $\rho$; taking limits proves upper
semicontinuity.

For lower semicontinuity when $p<1$, choose an optimal
$K=U\sqrt E$ for $\rho$. If its success is strictly larger than $p$,
it remains feasible in a neighbourhood. If the success equals $p$, use
$E_t=(1-t)E+tI$ and $K_t=U\sqrt{E_t}$. The success is
$p+t(1-p)>p$, while the filtered value tends to the optimal value as
$t\downarrow0$, by continuity of the positive matrix square root.
Each $K_t$ remains feasible for all sufficiently nearby states. This
proves lower semicontinuity.

At $p=1$, positivity and
$\tr[(I-K^\dagger K)\rho_A]=0$ imply that $K$ is an isometry on the
support of $\rho_A$. The joint state is supported there, and this isometry
extends to a qubit unitary. Thus every admissible output differs from
$\rho$ only by a local unitary, and
$\HF_A(\rho;1)=F_3(\rho)$ is continuous.
\end{proof}

\section{An elementary exact frontier for every pure two-qubit state}

The following specialization connects the general formula to the familiar
pure-state entanglement concentration problem \cite{vidal}. It is included
as a derived check and useful closed expression, not as a claim to have
introduced optimal pure-state concentration.

\begin{theorem}[Pure-state rate frontier]\label{thm:purefrontier}
Let the smaller squared Schmidt coefficient be $\varepsilon\in[0,1/2]$.
If $\varepsilon=0$, both one-sided frontiers equal one. If
$0<\varepsilon\leq1/2$, then for either side
\begin{equation}\label{eq:purefrontier}\boxed{
 \HF_W(\rho;p)^2=
 \begin{cases}
 3,&0<p\leq2\varepsilon,\\[2mm]
 1+\dfrac{8\varepsilon(p-\varepsilon)}{p^2},&2\varepsilon\leq p\leq1.
 \end{cases}}
\end{equation}
All values are attained by invertible filters when $\varepsilon>0$.
\end{theorem}
\begin{proof}
In Schmidt coordinates the state is
$\sqrt{1-\varepsilon}|00\rangle+\sqrt\varepsilon|11\rangle$.
Rescale a filter to have success exactly $p$. The unnormalized Bob
marginal $B$ obeys
$0\preceq B\preceq D=\diag(1-\varepsilon,\varepsilon)$ and $\tr B=p$.
Writing its smaller eigenvalue as $s$, we have
$s\leq\min\{\varepsilon,p/2\}$, and
$\det B=s(p-s)$ is maximized by this largest allowed $s$.
A pure two-qubit state's correlation singular values are
$1,C,C$, where $C=2\sqrt{\det(B/p)}$. Hence
$F_3^2=1+8\det B/p^2$, giving both upper bounds in
\eqref{eq:purefrontier}.

For $p\geq2\varepsilon$, the positive invertible contraction
\[
 K=\diag\left(\sqrt{\frac{p-\varepsilon}{1-\varepsilon}},1\right)
\]
attains the bound. At $p=2\varepsilon$ it produces a Bell state.
Scaling that filter down gives the same Bell state with any smaller
positive probability. The product-state endpoint is immediate.
\end{proof}

For fixed $p>0$, \eqref{eq:purefrontier} tends to one as
$\varepsilon\downarrow0$, whereas the unrestricted squared optimum is
three for every $\varepsilon>0$. This gives an exact example of the
regularization mechanism in Theorem \ref{thm:fixedratecontinuous}.

\section{Convex geometry and the complete discontinuity locus}

\begin{theorem}[Quasiconvexity, but not convexity]\label{thm:quasiconvex}
For any two density matrices and $0\leq t\leq1$,
\begin{equation}\label{eq:quasiconvex}
 \HF_A(t\rho+(1-t)\sigma)
 \leq\max\{\HF_A(\rho),\HF_A(\sigma)\}.
\end{equation}
Every sublevel set is closed and convex. The function itself is not convex.
\end{theorem}
\begin{proof}
For a fixed invertible filter, the filtered mixture is a mixture of the
two filtered states with weights proportional to their positive success
probabilities. The map $F_3$ is convex because it is a norm of a linear
function of the state. This gives the maximum bound for that filter;
taking the supremum gives \eqref{eq:quasiconvex}.

More explicitly, for any fixed local filter $K$ let $N_K(\rho)$ be the
nine unnormalized correlation coordinates and let $P_K(\rho)$ be its
success probability. Both are real-linear in $\rho$. For $v\geq0$,
\[
 \{\rho:\HF_A(\rho)\leq v\}
 =\bigcap_K\{\rho:\norm{N_K(\rho)}_2\leq vP_K(\rho)\}.
\]
Each condition is a closed convex norm inequality; negative sublevels are empty. Singular filters may
also be included: a zero-success positive output has trace zero and is
the zero matrix, making its condition $0\leq0$.

For nonconvexity take
$\rho_1=|00\rangle\langle00|$ and
$\rho_2=|1\rangle\langle1|\otimes I/2$.
Their Alice-filtered values are respectively one and zero, but every
$t\rho_1+(1-t)\rho_2$, $0<t<1$, has supremum one by postselection on
Alice's $|0\rangle$. It cannot exceed one because it is separable.
Indeed a product correlation matrix has Frobenius norm at most one,
and convexity proves the same for mixtures of products; local filtering
preserves separability. Thus the mixture value is $1>t$.
\end{proof}

\begin{theorem}[Exact continuity locus]\label{thm:continuitylocus}
The unrestricted quantity $\HF_A$ is lower semicontinuous everywhere.
It is continuous at a state if and only if that state's Alice marginal
is mixed. For $\HF_B$, replace Alice by Bob.
\end{theorem}
\begin{proof}
For every fixed invertible $K$, the success denominator is bounded below
by the positive smallest eigenvalue of $K^\dagger K$, uniformly over all
states. Its filtered value is therefore continuous. A supremum of
continuous functions is lower semicontinuous.

At a state with $\norm a<1$, choose a neighbourhood where
$\norm a\leq r<1$. The fractional formula \eqref{eq:fraction} is jointly
continuous on that neighbourhood times the compact closed unit ball,
with denominator bounded below by $(1-r)^2$. Its maximum is continuous,
for example by uniform continuity and
$|\max f-\max g|\leq\max|f-g|$.

At any pure-Alice-marginal state, use local unitaries to write
$\rho_0=|0\rangle\langle0|\otimes\sigma_B$ with Bloch vector $b$.
If $\norm b<1$, the separable states
\[
 \rho_\varepsilon=(1-\varepsilon)|0\rangle\langle0|\otimes\sigma_B
 +\varepsilon|1\rangle\langle1|\otimes|\phi\rangle\langle\phi|
\]
converge to $\rho_0$ and have $\HF_A=1$, whereas
$\HF_A(\rho_0)=\norm b<1$. If $\norm b=1$, locally rotate to
$|00\rangle$ and approach it by the entangled pure Schmidt states of
Theorem \ref{thm:purefrontier}. Their unrestricted value is $\sqrt3$,
whereas the product state's value is one. Every pure-marginal point is
therefore a discontinuity point, completing the classification.
\end{proof}

\section{Quantitative optimality and numerical stability}

\begin{proposition}[Distance and conditioning from the exact gap]
\label{prop:stability}
Let $V=\HF_A^2$ and retain $A,x_0,\lambda$ from the main theorem.
If $A\succeq mI$ with $m>0$, then every feasible $x$ satisfies
\begin{equation}\label{eq:coercive}
 V-q(x)\geq m\norm{x-x_0}^2.
\end{equation}
Thus an $\epsilon$-optimal effect lies within
$\sqrt{\epsilon/m}$ of the unique optimal effect. If $A$ is singular
and nonzero, the analogous bound uses its smallest positive eigenvalue
and distance to $x_0+\ker A$. If $A=0$, that distance is identically
zero, so only the trivial distance bound applies.
If $\lambda>0$ and $V-q(x)\leq\epsilon$, then
\begin{equation}\label{eq:condition}
 1-\norm x^2\leq\frac\epsilon\lambda,
 \qquad\frac{1-\norm x}{1+\norm x}\leq\frac\epsilon\lambda.
\end{equation}
The last ratio is the smaller-to-larger eigenvalue ratio of the canonical
effect $I+x\cdot\sigma$.
\end{proposition}
\begin{proof}
Both statements follow from
$V-q(x)=\lambda(1-\norm x^2)+(x-x_0)^TA(x-x_0)$.
For singular $A$ use its spectral decomposition on $(\ker A)^\perp$.
For the eigenvalue ratio observe
$(1-r)/(1+r)=(1-r^2)/(1+r)^2\leq1-r^2$ when $0\leq r\leq1$.
\end{proof}

\begin{proposition}[Stability under coefficient perturbations]
For two real symmetric quadratic problems on the unit ball,
\[
 |V(\tau,H,g)-V(\widetilde\tau,\widetilde H,\widetilde g)|
 \leq|\tau-\widetilde\tau|+\norm{H-\widetilde H}_{\rm op}
          +2\norm{g-\widetilde g}.
\]
\end{proposition}
\begin{proof}
The same expression bounds the absolute difference of their objective
values at every point of the ball. Taking maxima in both directions proves
the result. This avoids differentiating a repeated or singular selected root.
\end{proof}

The continuity theorem also explains a critical implementation rule:
a nearly pure marginal must never be silently treated as exactly pure.
Arbitrarily weakly entangled pure states have squared unrestricted optimum
three, even when ordinary floating-point arithmetic rounds their marginal
purity to one. The supplied implementations conservatively reject such
ambiguous cases and request higher precision. Their behavior is tested
with explicit very small Schmidt weights.

\section{The complete interval of attainable values}

For mixed filtered marginals, not only the largest value but the entire
closed-filter range is algebraically determined. For any symmetric $H$
and vector $g$, let
\[
 \mathcal V(\tau,H,g)=\tau+\mu+g^T(\mu I-H)^\dagger g,
 \quad\mu=\max\left\{0,\rmax
   \begin{pmatrix}H&I\\gg^T&H\end{pmatrix}\right\}.
\]
\begin{corollary}[Exact attainable interval]
For the canonical quadratic, the complete set of squared $F_3$ values
under positive-success closed filters is
\[
 \left[-\mathcal V(-\tau,-H,-g),\ \mathcal V(\tau,H,g)\right].
\]
Its endpoints have the same explicit pseudoinverse/kernel classifications,
using $(\tau,H,g)$ for the upper endpoint and $(-\tau,-H,-g)$ for the lower.
Every value strictly between the endpoints is attained by an invertible
filter. An endpoint belongs to the invertible range exactly when its
corresponding optimal family meets the open ball.
\end{corollary}
\begin{proof}
Negating the quadratic gives its minimum by the same maximum theorem.
The closed ball is connected and compact, so its continuous real-valued
image is the interval between its minimum and maximum. The open ball is
connected and dense in the closed ball. Its image is an interval whose
closure has these same two endpoints, so it contains every intermediate
value. Endpoint attainment is exactly the already proved strict-ball
criterion. Taking square roots gives the corresponding $F_3$ interval.
\end{proof}
